\section{Instrucciones de utilización de las simulaciones} \label{sec:instrucciones}

% Esta sección debe mostrar los comandos necesarios para hacer funcionar la simulación en todos los casos que especifica el plan de pruebas.
% Hay que suponer que el diseño de un grupo puede ser utilizado por otro grupo o el profesor.
% Si los resultados no se pueden repetir porque no se conocen los comandos para hacer funcionar la simulación entonces es como si el diseño no funcionara del todo.
% Se recomienda crear un Makefile de modo que se pueda correr todas las pruebas del caso con un solo comando en Icarus Verilog y GTKwave.
Para realizar la simulación del programa y verificar la correcta ejecución de las pruebas realizadas para cada submódulo y de forma completa, se utilizaron \texttt{Makefile} que agrupan los comandos e instrucciones necesarias.
Contienen las reglas necesarias para compilar con \texttt{iverilog}, ejecutar la simulación, generar el archivo de ondas \texttt{VCD} y abrir \texttt{GTKWave} con una configuración de señales definida.

\subsection*{Requisitos de software}

\begin{itemize}
  \item \texttt{iverilog} (>12.0) (Icarus Verilog) para compilar el diseño estructural.
  \item \texttt{vvp} (intérprete de Icarus Verilog) (>12.0) para ejecutar la simulación.
  \item \texttt{gtkwave} para visualizar las señales simuladas.
  \item \texttt{make} para automatizar la compilación y simulación.
\end{itemize}

\subsection*{Ejecución}

Respecto a la compilación y ejecución de las simulaciones, se implementaron \texttt{Makefile} dentro del directorio de cada submódulo en la carpeta \texttt{src/}.
Estos son controlados por un \texttt{Makefile} maestro en la raíz del repositorio.
Por lo tanto, para la ejecución, se realiza desde la raíz del repositorio y basta con colocar el nombre del módulo que se desea simular después del comando \texttt{make}.
Si no se coloca nada, \texttt{make} simula el PCS completo automáticamente.

Entonces, para simular todo el PCS, utilice el siguiente comando:
\begin{minted}{shell}
make
#o
make pcs
\end{minted}

Para simular un único módulo, se realiza de la siguiente forma:
\begin{minted}{shell}
make [transmit|receive|sync]
\end{minted}

\subsection*{Eliminar archivos generados}

Utilice el comando a continuación para eliminar los archivos de salida del proceso de compilación y simulación, estos son los generados en todo el repositorio:

\begin{minted}{shell}
make clean
\end{minted}
